% !PRINTED 126
\weberpage{126}

%<<CONT
be. For if we set $t = m s + s_1$, where $s_1 < s$, then $c^{s_1}$
is contained in $Q$, and $s_1$ must therefore be $= 0$. Accordingly $s$ is a
divisor of $\gamma$. But now the elements

\[
Q + c\,Q + c^2\,Q + \cdots c^{s-1}\,Q
\]

form, by (1), a group contained in $P$, and indeed of degree $p^{\varkappa} s$,
and since $p^{\varkappa} s$ must then be a divisor of $n$, it follows that $s$
cannot be divisible by $p$. Now if the degree $\gamma$ of $c$
were a power of $p$, then $s$ would also have to be a power of $p$, which
by the foregoing is not possible.

The degree of $c$ can therefore certainly not be a power of
$p$; i.e., \gesperrt{apart from the elements $Q$ there are in $R$ no
elements whose degree is a power of $p$.}

If now $b$ denotes an element of $P$ that does not belong to $R$,
then by the definition of $R$ there is in the group

\[
Q' = b^{-1}\,Q\,b
\]

at least one element that does not occur in $Q$, and the
degree of this element must be a power of $p$, because the
degree of $Q'$ is equal to $p^{\varkappa}$, hence a power of $p$. This element
can therefore not occur in $R$, \gesperrt{and $Q'$ is accordingly not a
divisor of $R$.}

If we choose the elements $b_1, b_2, \ldots, b_{j-1}$ from $P$ in such a way
that we can set

\weq{10}{P = R + R\,b_1 + R\,b_2 + \cdots + R\,b_{j-1}} % [sic: Weber jumps from equation (2) to (10) in this section and subsequently cites (10).]

then the $j - 1$ groups:

\[
b_1^{-1}\,Q\,b_1,\ \ b_2^{-1}\,Q\,b_2,\ \ \ldots,\ \ b_{j-1}^{-1}\,Q\,b_{j-1},
\]

which are all of degree $p^{\varkappa}$, are distinct from $Q$. But they are
also distinct from one another; for if we had, e.g.,

\[
b_1^{-1}\,Q\,b_1 = b_2^{-1}\,Q\,b_2,
\]

it would follow that:

\[
b_2\,b_1^{-1}\,Q\,b_1\,b_2^{-1} = Q,
\]

i.e., $b_2\,b_1^{-1} = c$ would be contained in $R$, hence $b_2 = c\,b_1$ in $R\,b_1$, which
by (10) is not possible. We therefore certainly have $j$ distinct
conjugate groups of degree $p^{\varkappa}$:

\weq{11}{Q,\ \ b_1^{-1}\,Q\,b_1,\ \ b_2^{-1}\,Q\,b_2,\ \ \ldots,\ \ b_{j-1}^{-1}\,Q\,b_{j-1}.}

If $Q' = b^{-1}\,Q\,b$ is one of them, then $R' = b^{-1}\,R\,b$ contains
the divisor $Q'$, but apart from this no element whose degree
%CONT>>
